Přímé srovnání mediánu sjednaných mzdových nárůstů v~\ac{KS} s~inflací \ac{HICP} ukazuje, že ČR prošla v~průběhu sledovaného období dvěma výraznými epizodami reálného poklesu kupní síly: poprvé v~době globální finanční krize (2009--2010) a~podruhé a~výrazněji během inflační vlny 2021--2023.
V~obou případech kolektivní smlouvy sjednané na předchozí rok neposkytovaly dostatečné krytí --- a~to navzdory tomu, že inflační tlak byl ekonomickými analytiky předvídatelný.
Tento selhávající mechanismus ochrany reálných mezd je přímým důsledkem decentralizované a~krátkodobé struktury \ac{KS} v~ČR: smlouvy jsou uzavírány převážně na podnikové úrovni, na jeden rok, bez inflačních doložek.
Sektorové smlouvy s~víceletou platností a~automatickými inflačními klausulemi, jak je standardní v~Belgii nebo Luxembursku, by podobné eroze reálných mezd systémově eliminovaly.

Přímé srovnání mediánu sjednaných mzdových nárůstů v~\ac{KS} s~inflací \ac{HICP} ukazuje, že ČR prošla v~průběhu sledovaného období dvěma výraznými epizodami reálného poklesu kupní síly: poprvé v~době globální finanční krize (2009--2010) a~podruhé a~výrazněji během inflační vlny 2021--2023.
V~obou případech kolektivní smlouvy sjednané na předchozí rok neposkytovaly dostatečné krytí --- a~to navzdory tomu, že inflační tlak byl ekonomickými analytiky předvídatelný.
Tento selhávající mechanismus ochrany reálných mezd je přímým důsledkem decentralizované a~krátkodobé struktury \ac{KS} v~ČR: smlouvy jsou uzavírány převážně na podnikové úrovni, na jeden rok, bez inflačních doložek.
Sektorové smlouvy s~víceletou platností a~automatickými inflačními klausulemi, jak je standardní v~Belgii nebo Luxembursku, by podobné eroze reálných mezd systémově eliminovaly.

Přímé srovnání mediánu sjednaných mzdových nárůstů v~\ac{KS} s~inflací \ac{HICP} ukazuje, že ČR prošla v~průběhu sledovaného období dvěma výraznými epizodami reálného poklesu kupní síly: poprvé v~době globální finanční krize (2009--2010) a~podruhé a~výrazněji během inflační vlny 2021--2023.
V~obou případech kolektivní smlouvy sjednané na předchozí rok neposkytovaly dostatečné krytí --- a~to navzdory tomu, že inflační tlak byl ekonomickými analytiky předvídatelný.
Tento selhávající mechanismus ochrany reálných mezd je přímým důsledkem decentralizované a~krátkodobé struktury \ac{KS} v~ČR: smlouvy jsou uzavírány převážně na podnikové úrovni, na jeden rok, bez inflačních doložek.
Sektorové smlouvy s~víceletou platností a~automatickými inflačními klausulemi, jak je standardní v~Belgii nebo Luxembursku, by podobné eroze reálných mezd systémově eliminovaly.

Přímé srovnání mediánu sjednaných mzdových nárůstů v~\ac{KS} s~inflací \ac{HICP} ukazuje, že ČR prošla v~průběhu sledovaného období dvěma výraznými epizodami reálného poklesu kupní síly: poprvé v~době globální finanční krize (2009--2010) a~podruhé a~výrazněji během inflační vlny 2021--2023.
V~obou případech kolektivní smlouvy sjednané na předchozí rok neposkytovaly dostatečné krytí --- a~to navzdory tomu, že inflační tlak byl ekonomickými analytiky předvídatelný.
Tento selhávající mechanismus ochrany reálných mezd je přímým důsledkem decentralizované a~krátkodobé struktury \ac{KS} v~ČR: smlouvy jsou uzavírány převážně na podnikové úrovni, na jeden rok, bez inflačních doložek.
Sektorové smlouvy s~víceletou platností a~automatickými inflačními klausulemi, jak je standardní v~Belgii nebo Luxembursku, by podobné eroze reálných mezd systémově eliminovaly.

\input{texparts/python/ipp_neg_vs_inflation}



